\documentclass[12pt, a4paper]{report}

% Packages
\usepackage[utf8]{inputenc}
\usepackage[T1]{fontenc}
\usepackage{geometry}
\usepackage{titlesec}
\usepackage{booktabs}
\usepackage{array}
\usepackage{enumitem}
\usepackage{hyperref}
\usepackage{listings}
\usepackage{xcolor}
\usepackage{tcolorbox}
\usepackage{tikz}
\usepackage{amsmath}
\usepackage{amssymb}
\usepackage{fancyhdr}
\usepackage{graphicx}
\usepackage{caption}
\usepackage{float}

% Geometry setup
\geometry{margin=1in, headheight=15pt}

% Custom colors
\definecolor{codegreen}{rgb}{0,0.6,0}
\definecolor{codegray}{rgb}{0.5,0.5,0.5}
\definecolor{codepurple}{rgb}{0.58,0,0.82}
\definecolor{backcolour}{rgb}{0.95,0.95,0.92}
\definecolor{primaryblue}{RGB}{0, 102, 204}
\definecolor{secondarygreen}{RGB}{0, 153, 76}
\definecolor{accentorange}{RGB}{255, 153, 0}

% Code listing style
\lstset{
    backgroundcolor=\color{backcolour},
    commentstyle=\color{codegreen},
    keywordstyle=\color{magenta},
    numberstyle=\tiny\color{codegray},
    stringstyle=\color{codepurple},
    basicstyle=\ttfamily\footnotesize,
    breakatwhitespace=false,
    breaklines=true,
    captionpos=b,
    keepspaces=true,
    numbers=left,
    numbersep=5pt,
    showspaces=false,
    showstringspaces=false,
    showtabs=false,
    tabsize=4,
    frame=single,
    language=Python
}

% Hyperref setup
\hypersetup{
    colorlinks=true,
    linkcolor=primaryblue,
    filecolor=magenta,      
    urlcolor=primaryblue,
    citecolor=secondarygreen,
}

% Section formatting
\titleformat{\chapter}[display]
  {\normalfont\huge\bfseries\color{primaryblue}}
  {\chaptertitlename\ \thechapter}{20pt}{\Huge}
\titleformat{\section}
  {\normalfont\Large\bfseries\color{primaryblue}}
  {\thesection}{10pt}{}
\titleformat{\subsection}
  {\normalfont\large\bfseries\color{secondarygreen}}
  {\thesubsection}{10pt}{}

% Custom commands
\newcommand{\techstack}[1]{\texttt{#1}}
\newcommand{\cost}[1]{\textbf{\color{secondarygreen}$#1}}
\newcommand{\warning}[1]{\textbf{\color{red}#1}}
\newcommand{\note}[1]{\begin{tcolorbox}[colback=blue!5!white,colframe=blue!75!black,title=Note]#1\end{tcolorbox}}
\newcommand{\important}[1]{\begin{tcolorbox}[colback=red!5!white,colframe=red!75!black,title=Important]#1\end{tcolorbox}}

% Title information
\title{\textbf{\Huge AI Middleman Contact Matching System}\\[1em]
\Large Comprehensive Technical Plan}
\author{System Architecture Document}
\date{\today}

\begin{document}

% Title Page
\begin{titlepage}
    \centering
    \vspace*{2cm}
    
    {\Huge\bfseries\color{primaryblue}AI Middleman}\\[0.5em]
    {\LARGE\bfseries Contact Matching System}\\[2cm]
    
    {\large Comprehensive Technical Plan}\\[1cm]
    
    \vspace{1cm}
    
    \begin{tabular}{|l|l|}
        \hline
        \textbf{Document Version} & 2.0 \\
        \hline
        \textbf{Date} & \today \\
        \hline
        \textbf{Status} & Planning Phase \\
        \hline
    \end{tabular}
    
    \vfill
    
    {\large\itshape "Automating Business Connections through Intelligent Matching"}
    
    \vspace{2cm}
\end{titlepage}

% Table of Contents
\tableofcontents
\newpage

% Chapter 1: Executive Summary
\chapter{Executive Summary}

\section{Project Overview}

This document presents a comprehensive technical plan for building a simplified, high-efficiency \textbf{AI-powered contact matching system} designed for a business middleman who maintains a network of 200--400 professional contacts. The system automates the process of matching incoming connection requests (received via WhatsApp) with appropriate contacts from the middleman's database using a streamlined agentic approach.

\section{Key Requirements}

\begin{itemize}
    \item \textbf{Contact Database:} 200--400 contacts with fields for name, contact number, email, role, company, context, and LinkedIn URL
    \item \textbf{Message Volume:} 50--250 WhatsApp messages per day requiring contact matching
    \item \textbf{Budget:} Free to low-cost (\cost{< \$5/month})
    \item \textbf{Latency:} Near real-time (< 3--5 seconds)
    \item \textbf{Deployment:} Free VPS or cloud infrastructure (Oracle Cloud Free Tier)
    \item \textbf{Primary Interface:} Mobile WhatsApp, with optional PC dashboard support
\end{itemize}

\section{Recommended Approach: Simple Agentic AI Architecture}

After analyzing several architectures and taking into account the small contact database size (200--400 entries), we recommend a \textbf{Simple Agentic AI Architecture} consisting of:

\begin{enumerate}
    \item \textbf{Stage 1: Keyword Filtering} to run locally and instantly, pruning the 400 contacts down to 30--50 potential candidates. This is completely free and eliminates the overhead and complexity of managing vector embeddings and local sentence transformer models.
    \item \textbf{Stage 2: LLM Agent Evaluation} via OpenRouter's free tier (Llama 3.1 8B) or a local Ollama fallback. The agent processes the 30--50 candidates within a single prompt, assessing suitability and outputting structured JSON containing confidence scores (0.0--1.0) and match reasoning.
    \item \textbf{No-Match Handling \& Clarification:} If the best match confidence is low (< 0.4), the system prompts the agent to ask clarifying questions back to the user on WhatsApp to gather additional requirements.
\end{enumerate}

This approach provides high accuracy, low latency, and absolute simplicity with minimal hosting requirements.

\section{Key Metrics}

\begin{center}
\begin{tabular}{|l|l|}
\hline
\textbf{Metric} & \textbf{Target} \\
\hline
Monthly Operating Cost & \cost{\$0-5} \\
\hline
Response Latency & 2--3 seconds \\
\hline
Match Accuracy & High \\
\hline
System Uptime & > 99\% \\
\hline
\end{tabular}
\end{center}

% Chapter 2: System Architecture
\chapter{System Architecture}

\section{High-Level Architecture}

\begin{figure}[H]
\centering
\begin{tikzpicture}[
    node distance=1.5cm and 1.5cm,
    box/.style={rectangle, draw, rounded corners, minimum width=2.5cm, minimum height=1.0cm, text centered, font=\small, fill=blue!10},
    db/.style={rectangle, draw, rounded corners, minimum width=2.5cm, minimum height=1.2cm, font=\small, fill=green!10},
    arrow/.style={->, >=stealth, thick}
]
    % WhatsApp
    \node[box, fill=green!20] (whatsapp) {WhatsApp\\Business API};
    
    % Webhook
    \node[box, right=of whatsapp] (webhook) {Webhook\\Receiver};
    
    % Message Processor
    \node[box, right=of webhook] (processor) {Message\\Processor};
    
    % Stage 1: Keyword Filter
    \node[box, fill=orange!10, below=of processor] (keyword) {Stage 1: Keyword\\Filter (30-50 cand.)};
    
    % Database
    \node[db, left=of keyword] (database) {PostgreSQL\\Database};
    
    % Stage 2: LLM Agent
    \node[box, fill=orange!20, below=of keyword] (agent) {Stage 2: LLM Agent\\(OpenRouter Llama)};
    
    % Response Formatter
    \node[box, fill=purple!15, left=of agent] (formatter) {Response\\Formatter};
    
    % Arrows
    \draw[arrow] (whatsapp) -- (webhook);
    \draw[arrow] (webhook) -- (processor);
    \draw[arrow] (processor) -- (keyword);
    \draw[arrow] (database) -- (keyword);
    \draw[arrow] (keyword) -- (agent);
    \draw[arrow] (agent) -- (formatter);
    \draw[arrow] (formatter) -- (database);
    \draw[arrow] (formatter) |- (whatsapp);
\end{tikzpicture}
\caption{High-Level System Architecture}
\label{fig:architecture}
\end{figure}

\section{Component Description}

\subsection{WhatsApp Business API Integration}
The system receives incoming messages through Meta's WhatsApp Business Cloud API via webhooks. This is the official, legal method for programmatic WhatsApp access.

\subsection{Webhook Receiver (FastAPI)}
A Python FastAPI server that:
\begin{itemize}
    \item Receives webhook payloads from WhatsApp
    \item Validates message authenticity
    \item Routes messages to the processing pipeline
    \item Returns generated matches or clarifying questions to users
\end{itemize}

\subsection{Message Processor}
Processes incoming messages by:
\begin{itemize}
    \item Parsing query text and extracting search tokens
    \item Saving raw inbound messages to the database
    \item Initiating the matching pipeline
\end{itemize}

\subsection{Matching Engine (Two-Stage Pipeline)}
The core logic of the system, simplified down to two stages to accommodate the small size of the database:
\begin{enumerate}
    \item \textbf{Stage 1: Keyword Filter:} A fast, local keyword filter queries the PostgreSQL database using tokens extracted from the request. It filters down the 400 contacts to approximately 30--50 candidates.
    \item \textbf{Stage 2: LLM Agent:} The candidate details are compiled into a prompt sent to Llama 3.1 8B (via OpenRouter free tier). The LLM acts as an agent, evaluating candidates, computing confidence scores, generating match reasoning, and outputting a structured JSON document.
\end{enumerate}

\subsection{Database (PostgreSQL)}
A relational database storing contact records, message history, match history, and agent reasoning. Unlike the hybrid plan, a vector index (\techstack{pgvector}) is not required, minimizing setup effort and system resources.

\subsection{Response Formatter \& Decision Branch}
Takes the structured output from the LLM agent and applies the decision rules:
\begin{itemize}
    \item \textbf{Good Match (Confidence > 0.7):} Returns the top 3--5 matches with details and reasoning.
    \item \textbf{Weak Match (Confidence 0.4--0.7):} Returns matches with a caveat indicating a partial/weak match.
    \item \textbf{No Match (Confidence < 0.4):} Asks a clarifying question to the WhatsApp user (e.g., "Found consultants in Johannesburg - would that work?").
\end{itemize}

% Chapter 3: Approach Comparison
\chapter{Approach Comparison}

\section{Comparison of Architectures}

For a database of 200--400 contacts, we compared the original five-stage Hybrid approach with the newly adopted two-stage Agentic AI approach.

\begin{table}[H]
\centering
\renewcommand{\arraystretch}{1.5}
\begin{tabular}{|l|c|c|}
\hline
\textbf{Criteria} & \textbf{Hybrid Architecture} & \textbf{Agentic AI Architecture (Selected)} \\
\hline
\textbf{Cost/Month} & \cost{\$2-9} & \cost{\$0-5} \\
\hline
\textbf{Latency} & 1--3 seconds & 2--3 seconds \\
\hline
\textbf{Accuracy} & High & \textbf{High} \\
\hline
\textbf{Embeddings Required} & Yes (pgvector, local BGE) & \textbf{No} \\
\hline
\textbf{Vector Search} & Core component & \textbf{Removed} \\
\hline
\textbf{LLM Input Candidates} & Top 10 (pre-filtered by cross-encoder) & Top 30--50 (pre-filtered by keywords) \\
\hline
\textbf{LLM Provider} & Groq API & OpenRouter Free Tier (Llama 3.1 8B) \\
\hline
\textbf{Local Fallback} & Cross-encoder only (degraded) & Ollama local models \\
\hline
\textbf{Complexity} & Medium-High (Multiple ML models) & \textbf{Low (Simple Python + API)} \\
\hline
\end{tabular}
\caption{Comparison Matrix: Hybrid vs. Agentic AI}
\label{tab:comparison}
\end{table}

\section{Why the Simple Agentic AI Approach Wins}

The transition from the Hybrid model to a simpler Agentic AI approach is driven by the following factors:

\begin{itemize}
    \item \textbf{No Embedding Management:} We do not need to generate, store, index, or update embedding vectors, which completely eliminates the dependency on \techstack{pgvector} and local embedding/cross-encoder inference pipelines.
    \item \textbf{Low Resource Footprint:} Running SentenceTransformers and Cross-Encoders locally on a free-tier VPS (like Oracle Cloud AMD or low-memory ARM instances) can lead to memory pressure and instance termination. The Agentic AI approach offloads reasoning to an external API (OpenRouter) and uses a simple, fast keyword filter locally.
    \item \textbf{Flexible and Context-Rich LLM Context:} With a total database size of 200--400 contacts, a keyword filter can prune the list to 30--50 potential candidates. Modern LLMs like Llama 3.1 8B have a large context window (128k tokens) and can easily analyze 30--50 candidates along with their full database context in a single prompt.
    \item \textbf{Handling of Ambiguity (Clarifications):} By using a dedicated LLM Agent in Stage 2, the system can dynamically formulate clarifying questions when matching confidence is low, rather than silently failing or returning irrelevant results.
\end{itemize}

% Chapter 4: Detailed Component Design
\chapter{Detailed Component Design}
\label{ch:component_design}

\section{WhatsApp Integration}

\subsection{WhatsApp Business Cloud API}

\begin{tcolorbox}[colback=green!5!white,colframe=green!75!black,title=WhatsApp API Details]
\textbf{Free Tier:} 1,000 conversation sessions/month\\
\textbf{Cost After Free:} \cost{\$0.005}/conversation\\
\textbf{Estimated Monthly Cost:} \cost{\$1-3} (for 50--250 messages/day)
\end{tcolorbox}

\subsection{Webhook Configuration}

\begin{lstlisting}[language=JSON, caption=Webhook Payload Structure]
{
    "object": "whatsapp_business_account",
    "entry": [{
        "id": "WHATSAPP_BUSINESS_ACCOUNT_ID",
        "changes": [{
            "value": {
                "messaging_product": "whatsapp",
                "metadata": {
                    "phone_number_id": "PHONE_NUMBER_ID",
                    "display_phone_number": "+1234567890"
                },
                "contacts": [...],
                "messages": [{
                    "id": "MESSAGE_ID",
                    "from": "SENDER_NUMBER",
                    "timestamp": "1234567890",
                    "type": "text",
                    "text": {
                        "body": "Looking for a marketing consultant in Johannesburg"
                    }
                }]
            },
            "field": "messages"
        }]
    }]
}
\end{lstlisting}

\subsection{Webhook Receiver Implementation}

\begin{lstlisting}[language=Python, caption=FastAPI Webhook Receiver]
from fastapi import FastAPI, Request, HTTPException
import hashlib, hmac
import os

app = FastAPI()
APP_SECRET = os.getenv("WHATSAPP_APP_SECRET")

def verify_signature(payload: bytes, signature: str) -> bool:
    expected = "sha256=" + hmac.new(
        APP_SECRET.encode(), payload, hashlib.sha256
    ).hexdigest()
    return hmac.compare_digest(expected, signature)

@app.post("/webhook/whatsapp")
async def whatsapp_webhook(request: Request):
    payload = await request.body()
    signature = request.headers.get("X-Hub-Signature-256", "")
    if not verify_signature(payload, signature):
        raise HTTPException(status_code=401, detail="Invalid signature")
    
    data = await request.json()
    if data.get("object") != "whatsapp_business_account":
        return {"status": "ignored"}
    
    for entry in data["entry"]:
        for change in entry["changes"]:
            if change["field"] == "messages":
                await process_message(change["value"])
    
    return {"status": "received"}

async def process_message(value: dict):
    messages = value.get("messages", [])
    if not messages:
        return
    
    message = messages[0]
    if message.get("type") != "text":
        return
    
    sender = message["from"]
    text = message["text"]["body"]
    
    # Process through matching pipeline (Stage 1 & 2)
    matches_response = await matching_engine.match_contact(text)
    
    # Format response and send back via WhatsApp
    response_text = format_response(matches_response)
    await whatsapp_client.send_message(
        to=sender,
        text=response_text
    )
\end{lstlisting}

\section{Matching Engine Design}

\subsection{Stage 1: Keyword Filter}

The local keyword filter matches query tokens against multiple fields (\texttt{name}, \texttt{role}, \texttt{company}, \texttt{context}) in the database. This acts as a coarse filter to prune the database of 400 contacts down to a candidate set of 30--50 relevant contacts.

\begin{lstlisting}[language=Python, caption=Keyword Filter Implementation]
import re
from typing import List, Dict

class KeywordFilter:
    def __init__(self, db_pool):
        self.db_pool = db_pool
    
    async def filter_candidates(self, query: str) -> List[Dict]:
        # Extract alphanumeric search tokens of length 3 or more
        tokens = [t.lower() for t in re.findall(r'\b\w{3,}\b', query)]
        if not tokens:
            return []
        
        # Build SQL pattern matches using a regular expression union
        regex_pattern = "|".join(tokens)
        
        query_sql = """
            SELECT id, name, role, company, context, contact_number, email, linkedin_url
            FROM contacts
            WHERE LOWER(name) ~ $1 
               OR LOWER(role) ~ $1 
               OR LOWER(company) ~ $1 
               OR LOWER(context) ~ $1
            LIMIT 50
        """
        
        async with self.db_pool.acquire() as conn:
            results = await conn.fetch(query_sql, regex_pattern)
            
        return [dict(row) for row in results]
\end{lstlisting}

\subsection{Stage 2: LLM Agent}

The LLM Agent evaluates the query against the 30--50 filtered candidates. The agent runs on Llama 3.1 8B via OpenRouter's free tier. 

\begin{lstlisting}[language=Python, caption=LLM Agent Implementation (OpenRouter)]
import httpx
import os
import json
from typing import Dict, Any

class LLMAgent:
    def __init__(self):
        self.api_key = os.getenv("OPENROUTER_API_KEY")
        self.api_url = "https://openrouter.ai/api/v1/chat/completions"
        self.model = "meta-llama/llama-3.1-8b-instruct:free"
        self.ollama_url = os.getenv("OLLAMA_URL", "http://localhost:11434/api/generate")

    async def evaluate_matches(self, query: str, candidates: list) -> Dict[str, Any]:
        if not candidates:
            return {
                "analysis": "No candidates to evaluate",
                "matches": [],
                "match_quality": "none",
                "clarification_question": "I couldn't find any relevant contacts. Could you tell me more about who you are looking for?"
            }
            
        prompt = self._build_prompt(query, candidates)
        
        try:
            async with httpx.AsyncClient() as client:
                response = await client.post(
                    self.api_url,
                    headers={
                        "Authorization": f"Bearer {self.api_key}",
                        "HTTP-Referer": "https://aimiddleman.org",
                        "Content-Type": "application/json"
                    },
                    json={
                        "model": self.model,
                        "messages": [{"role": "user", "content": prompt}],
                        "response_format": {"type": "json_object"},
                        "temperature": 0.1
                    },
                    timeout=15.0
                )
            
            if response.status_code == 200:
                result = response.json()
                content = result['choices'][0]['message']['content']
                return json.loads(content)
            else:
                raise httpx.HTTPStatusError("API Error", request=response.request, response=response)
                
        except Exception:
            # Fallback to local Ollama instance if OpenRouter rate limits or fails
            return await self._fallback_ollama(prompt)

    async def _fallback_ollama(self, prompt: str) -> Dict[str, Any]:
        try:
            async with httpx.AsyncClient() as client:
                response = await client.post(
                    self.ollama_url,
                    json={
                        "model": "llama3.1:8b",
                        "prompt": prompt + "\nRespond with raw JSON only matching the schema.",
                        "format": "json",
                        "stream": False
                    },
                    timeout=20.0
                )
                if response.status_code == 200:
                    return json.loads(response.json()["response"])
        except Exception:
            pass
            
        # Emergency hard fallback
        return {
            "analysis": "LLM failed completely",
            "matches": [],
            "match_quality": "none",
            "clarification_question": "Sorry, I am experiencing temporary system difficulties. Please try again soon."
        }

    def _build_prompt(self, query: str, candidates: list) -> str:
        candidate_list = "\n".join(
            f"ID: {c['id']} | Name: {c['name']} | Role: {c['role']} | Company: {c.get('company','')} | Context: {c.get('context','')}"
            for c in candidates
        )
        return f"""Evaluate these candidates for the query: "{query}"
        
Candidates:
{candidate_list}

Respond ONLY with a JSON object following this exact schema:
{{
  "analysis": "Brief analysis of query requirements",
  "matches": [
    {{
      "contact_id": 123,
      "name": "Candidate Name",
      "role": "Role Title",
      "confidence": 0.85,
      "reasoning": "Detailed reason why this contact fits this query"
    }}
  ],
  "match_quality": "good" | "weak" | "none",
  "clarification_question": "Optional question if confidence is low"
}}"""
\end{lstlisting}

\subsection{Agent Output Format}

The agent outputs a structured JSON document representing its evaluation:

\begin{lstlisting}[language=JSON, caption=Agent Output JSON Schema]
{
  "analysis": "Brief analysis of what user needs",
  "matches": [
    {
      "contact_id": 123,
      "name": "John Smith",
      "role": "Marketing Consultant",
      "confidence": 0.85,
      "reasoning": "Why this is a good match..."
    }
  ],
  "match_quality": "good" | "weak" | "none",
  "clarification_question": "Optional question if needed"
}
\end{lstlisting}

\subsection{Response Formatter & Decision Branch}

The output is processed by the application routing code to determine the formatting of the WhatsApp message.

\begin{lstlisting}[language=Python, caption=Response Formatter Implementation]
from typing import Dict, Any

def format_response(agent_output: Dict[str, Any]) -> str:
    quality = agent_output.get("match_quality", "none")
    matches = agent_output.get("matches", [])
    clarification = agent_output.get("clarification_question", "")
    
    if quality == "good" and matches:
        text = "Here are the top matches I found:\n\n"
        for idx, m in enumerate(matches[:5]):
            text += f"{idx+1}. *{m['name']}*\n"
            text += f"   Role: {m['role']}\n"
            text += f"   Reasoning: {m['reasoning']}\n\n"
        return text.strip()
        
    elif quality == "weak" and matches:
        text = "⚠️ *Please note*: I found some partial matches that might fit, but they may not meet all criteria:\n\n"
        for idx, m in enumerate(matches[:3]):
            text += f"{idx+1}. *{m['name']}* ({m['role']})\n"
            text += f"   Reasoning: {m['reasoning']}\n\n"
        return text.strip()
        
    else:
        if clarification:
            return clarification
        return "I could not find any contacts matching your request. Could you please specify more details or keywords?"
\end{lstlisting}

% Chapter 5: Database Design
\chapter{Database Design}

\section{Schema Overview}

The database is built on standard PostgreSQL. Since the contact network is small (200--400 entries), we do not need embedding vectors or specialized index extensions like \techstack{pgvector}. Relational indexing is applied to query performance fields.

\begin{lstlisting}[language=SQL, caption=Database Schema]
-- Contacts table
CREATE TABLE contacts (
    id SERIAL PRIMARY KEY,
    name VARCHAR(255) NOT NULL,
    contact_number VARCHAR(20),
    email VARCHAR(255),
    role VARCHAR(255),
    company VARCHAR(255),
    context TEXT,
    linkedin_url VARCHAR(500),
    created_at TIMESTAMP DEFAULT CURRENT_TIMESTAMP,
    updated_at TIMESTAMP DEFAULT CURRENT_TIMESTAMP
);

-- Messages table
CREATE TABLE messages (
    id SERIAL PRIMARY KEY,
    sender_number VARCHAR(20) NOT NULL,
    message_text TEXT NOT NULL,
    received_at TIMESTAMP DEFAULT CURRENT_TIMESTAMP,
    processed BOOLEAN DEFAULT FALSE
);

-- Match history table
CREATE TABLE match_history (
    id SERIAL PRIMARY KEY,
    message_id INTEGER REFERENCES messages(id),
    contact_id INTEGER REFERENCES contacts(id),
    rank_position INTEGER NOT NULL,
    confidence FLOAT NOT NULL,
    reasoning TEXT,
    feedback VARCHAR(20),  -- 'accepted', 'rejected', 'pending'
    created_at TIMESTAMP DEFAULT CURRENT_TIMESTAMP
);

-- Indexes for performance
CREATE INDEX idx_contacts_name ON contacts(name);
CREATE INDEX idx_contacts_role ON contacts(role);
CREATE INDEX idx_messages_sender ON messages(sender_number);
\end{lstlisting}

% Chapter 6: Cost Analysis
\chapter{Cost Analysis}

\section{Monthly Cost Breakdown}

\begin{table}[H]
\centering
\renewcommand{\arraystretch}{1.4}
\begin{tabular}{|l|c|l|}
\hline
\textbf{Component} & \textbf{Cost/Month} & \textbf{Notes} \\
\hline
WhatsApp Business API & \cost{\$1-3} & 50--250 msg/day \\
\hline
LLM API (OpenRouter) & \cost{\$0} & Free Tier (Llama 3.1 8B) \\
\hline
VPS Hosting & \cost{\$0} & Oracle Cloud Free \\
\hline
PostgreSQL & \cost{\$0} & Self-hosted on VPS \\
\hline
Embedding / Cross-Encoder & \cost{\$0} & Not required \\
\hline
Domain + SSL & \cost{\$1} & Cheap domain \\
\hline
\textbf{Total} & \cost{\$1-4} & \textbf{Within budget} \\
\hline
\end{tabular}
\caption{Monthly Cost Breakdown}
\label{tab:costs}
\end{table}

\section{Free Tier Utilization}

\begin{itemize}
    \item \textbf{Oracle Cloud Free Tier:} 2 AMD VMs or ARM VMs (always free)
    \item \textbf{OpenRouter Free Tier:} Meta Llama 3.1 8B Instruct Free, providing rapid, cost-free evaluations
    \item \textbf{Ollama Local Fallback:} Self-hosted model on the VPS running local Llama 3.1 8B to recover from API failures or rate limits
    \item \textbf{PostgreSQL:} Standard open-source DB running locally on the VPS
\end{itemize}

% Chapter 7: Deployment Plan & Project Structure
\chapter{Deployment \& Project Structure}

\section{Project Structure}

The repository is structured to separate concern areas:

\begin{lstlisting}[caption=Project Directory Structure]
ai-middleman/
├── app/
│   ├── main.py              # FastAPI entrypoint
│   ├── database.py          # PostgreSQL connection
│   ├── models/              # Pydantic models
│   ├── routes/              # WhatsApp webhook, health
│   └── services/
│       ├── keyword_filter.py    # Stage 1: Local filter
│       ├── agent.py             # Stage 2: LLM Agent
│       ├── whatsapp_client.py   # Meta WhatsApp client
│       └── response_formatter.py # Decision branch formatter
├── tests/
├── migrations/
├── data/
└── docker-compose.yml
\end{lstlisting}

\section{Docker Compose Configuration}

The Docker setup builds the API server, spins up the database, and hosts a local Ollama instance for LLM fallback capability.

\begin{lstlisting}[language=bash, caption=docker-compose.yml]
version: '3.8'
services:
  db:
    image: postgres:15-alpine
    environment:
      POSTGRES_DB: aimiddleman
      POSTGRES_PASSWORD: ${DB_PASSWORD}
    volumes:
      - pgdata:/var/lib/postgresql/data
    ports:
      - "5432:5432"
  
  api:
    build: ./api
    ports:
      - "8000:8000"
    environment:
      DATABASE_URL: postgresql://postgres:${DB_PASSWORD}@db:5432/aimiddleman
      WHATSAPP_APP_SECRET: ${WHATSAPP_APP_SECRET}
      OPENROUTER_API_KEY: ${OPENROUTER_API_KEY}
      OLLAMA_URL: http://ollama:11434/api/generate
    depends_on:
      - db
      - ollama
  
  ollama:
    image: ollama/ollama:latest
    ports:
      - "11434:11434"
    volumes:
      - ollama:/root/.ollama

volumes:
  pgdata:
  ollama:
\end{lstlisting}

% Chapter 8: Implementation Roadmap
\chapter{Implementation Roadmap}

\section{Phase 1: Foundation (Day 1-2)}

\begin{itemize}
    \item Set up VPS and Docker environment
    \item Configure PostgreSQL database and create tables
    \item Implement basic FastAPI webhook receiver structures
    \item Build CSV data importer to load contact records
\end{itemize}

\section{Phase 2: Matching Engine (Day 3-4)}

\begin{itemize}
    \item Implement local token-based \techstack{KeywordFilter} (Stage 1)
    \item Build the \techstack{LLMAgent} prompt logic and integrate OpenRouter completions (Stage 2)
    \item Integrate local Ollama Llama 3.1 8B fallback path to catch rate limit exceptions
    \item Run unit and regression tests on search queries
\end{itemize}

\section{Phase 3: WhatsApp Integration (Day 5)}

\begin{itemize}
    \item Build Meta WhatsApp API client integration
    \item Implement \techstack{ResponseFormatter} for decision branching (good match, weak match, clarify)
    \item Complete end-to-end webhook receiver loops and verify WhatsApp replies
\end{itemize}

\section{Phase 4: Testing \& Polish (Day 6-7)}

\begin{itemize}
    \item Configure automated Caddy proxy with SSL certificate
    \item Set up monitoring logging systems
    \item Verify system resilience under artificial rate limit conditions
    \item Perform load and latency verification tests
\end{itemize}

% Chapter 9: Risk Assessment
\chapter{Risk Assessment}

\begin{table}[H]
\centering
\renewcommand{\arraystretch}{1.4}
\begin{tabular}{|l|c|l|}
\hline
\textbf{Risk} & \textbf{Impact} & \textbf{Mitigation} \\
\hline
WhatsApp API pricing shifts & High & Monitor conversation boundaries, restrict matching requests \\
\hline
OpenRouter rate limiting & Medium & Auto-fallback to local Ollama Llama 3.1 8B model \\
\hline
VPS downtime (Oracle) & Medium & Maintain current db backups; design fallback deploy (Fly.io/Render) \\
\hline
Clarification loops & Low & Set prompt rules instructing agent to limit questions \\
\hline
LLM schema failure & Low & Strict JSON parsing with try/except fallback blocks \\
\hline
\end{tabular}
\caption{Risk Assessment Matrix}
\end{table}

% Chapter 10: Conclusion
\chapter{Conclusion}

The two-stage Agentic AI architecture presents a highly cost-effective, straightforward, and performant implementation for a contact database of 200--400 professional profiles. By utilizing a simple keyword filtering stage before invoking Llama 3.1 8B via OpenRouter, the system achieves near real-time response times while operating on a minimal monthly budget (\cost{< \$5/month}).

The implementation removes pgvector indexing and embedding model layers, saving memory and CPU resources, while retaining high accuracy and conversational fallback capabilities to ensure excellent user experience.

\begin{tcolorbox}[colback=green!5!white,colframe=green!75!black,title=Next Steps]
\begin{enumerate}
    \item Obtain OpenRouter API key and register WhatsApp Business API account
    \item Provision Oracle Cloud VPS instance
    \item Set up repository according to the planned project structure
    \item Begin Phase 1 (Foundation - DB Setup and CSV Imports)
\end{enumerate}
\end{tcolorbox}

\end{document}
